\subsection{An explicit rational quotient of the cubic detector dessin}

\paragraph{D23.C14 (the marked genus-zero quotient and its cyclic pullback).}
Let $\xi:F\twoheadrightarrow\mathbb F_3^2\rtimes\operatorname{SL}_2(\mathbb F_3)$
be the marked cubic detector of D23.P10. Its degree-nine action has a
canonical extension to
\[
 \Gamma_+=\ker(C_2*C_3\longrightarrow C_2)
       =\langle u=r,v=srs\mid u^3=v^3=1\rangle.
\]
In the same affine origin as D23.P10, the images are
\[
 u=\left((2,2)^{\mathsf T},
       \begin{pmatrix}0&2\\1&2\end{pmatrix}\right),\qquad
 v=\left((1,1)^{\mathsf T},
       \begin{pmatrix}2&2\\1&0\end{pmatrix}\right).
\tag{D23.C14.1}
\]
Here $x=(vu)^{-1}$, and the branch-cycle types of $(x,v,u)$ are
\[
                   (621),\qquad(33111),\qquad(333).
\tag{D23.C14.2}
\]
The corresponding degree-nine rational Belyi map is
\[
 \Psi(z)=\frac{432z^6(z-1)(4z-1)^2}
                   {(19z^3-9z^2-3z+1)^3}.
\tag{D23.C14.3}
\]
Put
\[
 \begin{split}
 C(z)&=53z^3+15z^2-3z-1,\\
 Q(z)&=z^2-4z+1,\qquad H(z)=19z^3-9z^2-3z+1.
 \end{split}
\]
Its cyclic cubic pullback has the rational model
\[
                  w^3=C(z),
\tag{D23.C14.4}
\]
which is explicitly birational over $\mathbb Q$ to
$E:Y^2=X^3-4$. The formulas in both directions are
\[
 \begin{split}
 G&=-X^2-6X-3Y+10,\\
 z&=(2X+Y-2)/G,\qquad w=(X^2-2X+2Y+4)/G,\\
 X&=(4+16z-2w)/(1+z+w),\\
 Y&=6(w^2-19z^2-8z-1)/(1+z+w)^2.
 \end{split}
\tag{D23.C14.5}
\]
Let $n=(1+\sqrt{-3})/2$, $\bar n=1-n$, and let $\phi$ be the
normalized first-orbit degree-nine map in the retrieved primary BelyiDB data,
with $x_{\rm old}=9X/1600$ and $y_{\rm old}=27Y/64000$.
Then its compact expression and the precise composition identity are
\[
 \phi=\bar n\,\frac{H+wQ}{H-nwQ},\qquad
 \boxed{\ \Psi(z)=1-\left(\frac{\phi-\bar n}{\phi-n}\right)^3.\ }
\tag{D23.C14.6}
\]
Thus the square formed by $E\xrightarrow{\phi}\mathbb P^1_\lambda$,
$E\xrightarrow{z}\mathbb P^1_z$, $\mathbb P^1_z\xrightarrow{\Psi}\mathbb P^1_b$,
and $b=1-((\lambda-\bar n)/(\lambda-n))^3$ commutes with the stated
marking, not just with the stated passports.

The map $E\to\mathbb P^1_z$ is cyclic of degree three over
$\mathbb Q(\sqrt{-3})$ and has three ramification points, the points
$w=0$. It is therefore the quotient by a geometric order-three
symmetry with fixed points. In particular, the fact that $u$ has no
fixed sheet in (D23.C14.2) does \emph{not} make this symmetry a
fixed-point-free translation: the fixed sheets of $v$ account for
its three fixed points.
More explicitly, $w\mapsto\zeta^2w$, with $\zeta=n^2$, is
\[
       \alpha(Q)=\zeta^2Q+(1-\sqrt{-3},2\sqrt{-3})
\]
in the group law on $E$. It cycles the sixfold ramification points
$O$, $(1-\sqrt{-3},2\sqrt{-3})$, and
$(1+\sqrt{-3},-2\sqrt{-3})$.

\begin{proof}
The subgroup presentation is the index-two Reidemeister--Schreier
presentation of $C_2*C_3$. Restrict D23.P3's affine representation to
$\Gamma_+$. Its linear matrices then lie in
$\operatorname{SL}_2(\mathbb F_3)$, so taking the real part of the
translation is a homomorphism. The origin shift $(0,2)^{\mathsf T}$
from D23.P10 gives (D23.C14.1), and direct multiplication gives its
marked $x,y$. Their action on the nine affine points has the cycles
(D23.C14.2). In the labeling $(a,b)\mapsto3a+b$, the simultaneous
conjugator to the BelyiDB rational branch pair $(\sigma_0,\sigma_1)$ is
\[
                     [4,0,5,2,1,6,3,8,7].
\]
The certificate checks the two literal conjugacy equations. The
rational database passport has size one, and its map, after the
source coordinate change $x_{\rm DB}=7z$, factors as (D23.C14.3).
Equivalently, the formula itself has the required branch data since
\[
                 \Psi(z)-1=\frac{Q(z)^3C(z)}{H(z)^3},
\tag{D23.C14.7}
\]
all indicated factors are separable and coprime, and the partitions
in (D23.C14.2) have total ramification $16=2\cdot9-2$.

For the orbifold topology, let $N$ be the inverse image of a point
stabilizer in $\Gamma_+$, and put $U=N\cap F$. There are three
cusps, from the cycles of $x$, no elliptic points above $u$, and
three order-three elliptic points above $v$. Since
$\chi(\Gamma_+)=-1/3$ and $[\Gamma_+:N]=9$, the signature equation is
\[
 -3=2-2g(N)-3-3(1-1/3),
\]
so the compactification of $N$ has genus zero. A conjugate of $v$
fixing a sheet maps nontrivially to $\Gamma_+/F\simeq C_3$.
Consequently $[N:U]=3$ and $U\lhd N$. The three elliptic points are
exactly the ramification points of the compact cyclic cover.
The cusps split into nine cusps because $x\in F$. Therefore
$\chi(U)=-9=2-2g(U)-9$, and its compactification has genus one.
This also proves the fixed-point assertion without choosing a
potentially non-normalizing lift of $r$.

For the exact algebraic map, write
$L=2X+Y-2$ and $N_0=X^2-2X+2Y+4$. Modulo $Y^2-X^3+4$, one has
\[
 N_0^3=53L^3+15L^2G-3LG^2-G^3.
\]
This proves (D23.C14.4) under (D23.C14.5). Substitution verifies both
inverse formulas and the equation of $E$; hence these are
birational maps of smooth projective curves. Since $C$ is squarefree,
with discriminant $-11664$, $w^3=C(z)$ gives precisely three branch
points and no ramification at infinity.
The automorphism $w\mapsto\zeta^2w$ acts by $\zeta^2$ on the regular
differential $dz/w^2$, and its image of $O$ is
$(1-\sqrt{-3},2\sqrt{-3})$ by the inverse formulas. These two facts
give the displayed group-law expression for $\alpha$ and its orbit.

The normalized database numerator has the form $P(X)Y+R(X)$ and
denominator $D(X)$, recorded in the root's exact normalization file.
With
\[
 H_h=19L^3-9L^2G-3LG^2+G^3,\qquad
 Q_h=L^2-4LG+G^2,
\]
the following polynomial identity modulo $Y^2-X^3+4$ fixes the sign
and all scalar factors:
\[
 (PY+R)(H_h-nN_0Q_h)
       =\bar n D(H_h+N_0Q_h).
\]
It proves the first expression in (D23.C14.6). Solving for the
cross-ratio gives
\[
              \frac{\phi-\bar n}{\phi-n}
                       =n\,\frac{wQ}{H}.
\]
Because $n^3=-1$ and $w^3=C$, its cube is $-CQ^3/H^3$.
Equation (D23.C14.7) proves the boxed composition identity, with its
orientation as written. An earlier direct cross-multiplication of
the degree-27 composition checked the same identity and rejected
its opposite orientation; the displayed lower-degree identities
provide a shorter exact replay certificate.

For a normalization check, let $P_0=(2,2)$ on $E$. Then
$2P_0=(5,-11)$ and
\[
 (\phi)_0=6(O)+2(2P_0)+(-P_0),\qquad
 z(O)=0,\quad z(2P_0)=1/4,\quad z(-P_0)=1.
\]
The point with coordinates $(5,11)$ is not a zero of $\phi$.
These signs follow directly from the normalized numerator, not
from an abstract identification of the elliptic curves.

The exact scripts are
\path{agents/characters/code/check_cubic_orbifold_topology.py} and
\path{agents/characters/code/check_belyi_rational_quotient.py};
the corresponding JSON certificates are in
\path{agents/characters/records/}.
The primary database records are Michael Musty's BelyiDB,
\path{belyi_db/9/9T23-[6,3,3]-621-33111-333-g0.m} and
\path{belyi_db/9/9T23-[6,6,6]-621-621-621-g1.m},
fields \texttt{BelyiDBPermutationTriple}, \texttt{BelyiDBCurve},
and \texttt{BelyiDBBelyiMap}, first Galois orbit for the genus-one
record. The repository is
\url{https://github.com/michaelmusty/BelyiDB}.
The retrieved genus-zero record and its MIT license are retained
in the certificate directory. The normalized genus-one coefficients
are in \path{records/belyi_9T23_normalized.txt}; the marked
permutation comparison is in \path{records/belyi_9T23_match.json}.
\end{proof}

\paragraph{D23.C14a (the precise five-point comparison left by this diagram).}
In the coordinate
\[
                 t=(\lambda-\bar n)/(\lambda-n),\qquad b=1-t^3,
\]
the source of the cyclic cover is
$\mathbb P^1\setminus(\{0,\infty\}\cup\mu_3)$.
Writing $\zeta=n^2$, the points
$t=\zeta^2,\zeta,1$ are respectively $\lambda=0,1,\infty$,
and $t=0,\infty$ are $\lambda=\bar n,n$.
Thus recovering the original marked $\lambda$ sphere fills
$t=0,\infty$. The standard Enriquez distribution map instead
retains $t=0,1,\infty$ and fills the other two roots of unity.
The required relation between these two forgetting maps is an
additional marked comparison to be established; it is not being
asserted as a consequence of the existence of the diagram.

Also, $\Psi$ is indecomposable as a rational cover. Indeed its
monodromy is the affine group $\mathbb F_3^2\rtimes
\operatorname{SL}_2(\mathbb F_3)$ in its primitive degree-nine
action: a block containing zero is a translation subgroup stable
under the irreducible linear stabilizer. There is no proper such
subgroup. Consequently $\Psi$ cannot itself be a nontrivial
composition of degree-three rational Kummer covers. This rules out
that direct shortcut to applying distribution, while leaving the
five-point comparison available.
